\chapter[1976 --- Blumberg et al.]{1976: Baruch S. Blumberg, D. Carleton Gajdusek}
\label{ch:1976_Blumberg}

\section*{Main Contribution}

\textit{Discovered new mechanisms for the origin and spread of infectious diseases, including slow virus infections and the role of prion-like agents.}\cite{nobel-1976,lecture-1976}

\section{Official Citation}

for their discoveries concerning new mechanisms for the origin and dissemination of infectious diseases

The 1976 Nobel Prize in Physiology or Medicine was awarded jointly to Baruch S. Blumberg and D. Carleton Gajdusek for their independent but complementary discoveries concerning new mechanisms for the origin and dissemination of infectious diseases. Blumberg's discovery of the hepatitis B virus and its role in liver disease, and Gajdusek's elucidation of the infectious nature of kuru---a mysterious degenerative brain disease among the Fore people of Papua New Guinea---each revealed previously unknown mechanisms by which pathogens could cause disease. Together, their work expanded the boundaries of infectious disease biology in directions that neither could have anticipated, and their discoveries continue to resonate in modern medicine.

\section{Background and Historical Context}

The study of infectious diseases is one of the oldest and most consequential branches of medicine. By the middle of the twentieth century, the germ theory of disease---the principle that many diseases are caused by microorganisms---was well established, and major advances in bacteriology, virology, and parasitology had led to the development of vaccines, antibiotics, and public health measures that dramatically reduced the burden of infectious disease in the industrialized world. Yet major questions remained unanswered. Many common diseases, including chronic liver disease, certain forms of dementia, and a variety of degenerative neurological conditions, could not be attributed to any known pathogen. The mechanisms by which some diseases were transmitted---and the agents responsible for them---remained obscure.

Baruch Stanley Blumberg was born in Brooklyn, New York, in 1925, to a family of Eastern European Jewish immigrants. He trained in medicine at Columbia University and later pursued graduate studies in biochemistry and genetics at Oxford University. In the late 1950s, Blumberg began a systematic study of genetic variation in human populations, collecting blood samples from diverse populations around the world. His interest was in understanding how genetic polymorphisms might influence susceptibility to disease---a question that was at the forefront of human genetics at the time.

D. Carleton Gajdusek was born in Yonkers, New York, in 1923, and trained in medicine at Harvard Medical School. A voracious reader and an adventurous traveler, Gajdusek was drawn to the study of infectious diseases in tropical and remote regions of the world. In the 1950s and 1960s, he became fascinated by a devastating neurological disease called kuru, which was decimating the Fore people of the Eastern Highlands of Papua New Guinea. Kuru, which caused progressive cerebellar ataxia, tremors, and invariably death within months of symptom onset, was one of the great medical mysteries of the twentieth century.

The historical context of both discoveries was shaped by the post-World War II expansion of global health research. Organizations such as the World Health Organization (WHO), the National Institutes of Health (NIH), and the Rockefeller Foundation were supporting research into infectious diseases in developing countries, and there was a growing recognition that many diseases that had been neglected by Western medicine could yield important insights into fundamental biology.

\section{The Discovery/Contribution}

\subsection{Baruch Blumberg and the Discovery of Hepatitis B Virus}

Blumberg's journey to the discovery of hepatitis B virus (HBV) began with his study of genetic polymorphisms in human blood proteins. In the late 1950s, Blumberg and his colleague Alfred Pratt, working at the Fox Chase Cancer Center in Philadelphia, began collecting blood samples from populations around the world---including indigenous peoples of the Amazon, sub-Saharan Africa, and the Pacific Islands. They used a new technique called agar gel immunoelectrophoresis to identify previously unknown variations in serum proteins.

In 1965, while testing serum samples from an Australian Aboriginal patient, Blumberg and his colleague Harvey Alter discovered a previously unknown antigen in the blood of this patient. The antigen was later designated ``Australia antigen'' (Au antigen), and it appeared to be associated with a specific disease state. Blumberg initially hypothesized that the Au antigen was a genetic variant of a serum protein, since it appeared to be inherited in some populations. However, further investigation revealed a far more complex and important story.

Through a series of studies conducted between 1966 and 1968, Blumberg and his colleagues demonstrated that the Australia antigen was not a host protein at all, but rather a component of a virus---the virus that causes serum hepatitis (now called hepatitis B). They showed that the Au antigen was present in the blood of patients with serum hepatitis, that it was associated with the virus particles (later called Dane particles), and that it could be used as a serological marker for hepatitis B infection.

The identification of HBV and its associated antigens (including what are now known as hepatitis B surface antigen, or HBsAg) had immediate and significant implications for public health. Before the discovery of HBV, there was no way to screen blood donations for hepatitis infection, and transfusion-associated hepatitis was a major cause of morbidity and mortality. Blumberg's discovery made it possible to develop screening tests for HBV in blood supplies, dramatically reducing the incidence of transfusion-associated hepatitis. Furthermore, the characterization of the virus and its antigens led directly to the development of the hepatitis B vaccine---the first vaccine to be developed against a human cancer virus, since chronic HBV infection is a major cause of hepatocellular carcinoma (liver cancer).

Blumberg's work also had broader implications for the understanding of the relationship between infectious agents and cancer. The demonstration that a virus could cause a chronic infection that, in turn, could lead to cancer, was a landmark in cancer research. It confirmed the hypothesis, first proposed by Peyton Rous and others, that viruses could play a role in human cancer, and it provided a model system for studying the molecular mechanisms by which chronic infection leads to malignant transformation.

\subsection{D. Carleton Gajdusek and the Discovery of Kuru as an Infectious Disease}

Gajdusek's discovery of the infectious nature of kuru is one of the most notable stories in the history of medicine. Kuru was first described by Western medicine in the 1950s, when Australian administrators and missionaries in the Eastern Highlands of Papua New Guinea reported a mysterious disease affecting the Fore people and neighboring tribes. The disease caused a progressive loss of coordination (ataxia), uncontrollable trembling (hence the name ``kuru,'' which in the Fore language means ``to shiver'' or ``to tremble with fear''), emotional lability, and ultimately death, usually within six to twelve months of symptom onset.

Gajdusek first visited Papua New Guinea in 1957, at the invitation of Australian patrol officers who were concerned about the devastating effects of kuru on the Fore population. He was immediately struck by the epidemiological pattern of the disease: kuru occurred predominantly among women and children of the Fore tribe, and it appeared to be transmitted within families. The disease was not contagious in the conventional sense---close physical contact with kuru patients did not appear to spread the disease---but it was clearly infectious in some manner.

Working with Australian colleague Vincent Zigas, Gajdusek conducted a series of epidemiological and pathological studies of kuru in the 1950s and 1960s. He observed that the disease was associated with a particular cultural practice among the Fore: funerary cannibalism. When a member of the tribe died, the body was ritually prepared and consumed by the women and children of the community. Gajdusek hypothesized that kuru was transmitted through this practice of cannibalism---specifically, through the consumption of brain tissue from infected individuals.

To test this hypothesis, Gajdusek and his colleagues conducted a series of transmission experiments. They inoculated chimpanzees with brain tissue from kuru patients and, after a prolonged incubation period of 18 to 36 months, the chimpanzees developed a neurological disease that closely resembled kuru. This was a landmark achievement: it demonstrated that kuru was transmissible from humans to primates, providing powerful evidence that the disease was caused by an infectious agent.

However, the nature of this infectious agent remained a significant mystery. Gajdusek showed that the agent was remarkably resistant to conditions that would destroy conventional bacteria and viruses. It was not inactivated by formaldehyde, ultraviolet radiation, boiling, or ionizing radiation---all of which would destroy nucleic acids (the genetic material of conventional pathogens). This led Gajdusek to propose, in the 1960s, that the kuru agent was a ``slow virus''---an agent that replicated extremely slowly and caused disease only after a prolonged incubation period.

The true nature of the kuru agent would not be revealed until the 1980s, when Stanley Prusiner (who would share the Nobel Prize in 1997) proposed and demonstrated that kuru, along with other transmissible spongiform encephalopathies (TSEs) such as Creutzfeldt-Jakob disease and scrapie, was caused not by a virus but by a proteinaceous infectious particle---a prion. Gajdusek's work on kuru thus opened the door to one of the most unexpected and paradigm-shifting discoveries in the history of biology: the existence of infectious proteins that can propagate without nucleic acids.

\section{Impact on Modern Medicine}

The discoveries of Blumberg and Gajdusek have had far-reaching consequences for modern medicine. Blumberg's identification of hepatitis B virus and its associated antigens led directly to the development of blood screening tests that have saved millions of lives by preventing transfusion-associated hepatitis. The hepatitis B vaccine, developed in the 1980s and now included in the routine childhood vaccination schedule in most countries, has dramatically reduced the incidence of chronic hepatitis B infection and, consequently, the incidence of hepatocellular carcinoma. The World Health Organization estimates that hepatitis B vaccination has prevented approximately 300 million new cases of chronic hepatitis B infection worldwide.

The hepatitis B vaccine is also notable as the first vaccine to be developed against a virus that causes cancer. This achievement validated the concept that vaccination could be used not only to prevent infectious diseases but also to prevent cancer---a principle that has since been applied to the development of the human papillomavirus (HPV) vaccine, which prevents cervical and other cancers caused by HPV.

Gajdusek's work on kuru, meanwhile, opened up an entirely new field of research: the study of prion diseases. The recognition that kuru, Creutzfeldt-Jakob disease, Gerstmann-Str\"aussler-Scheinker syndrome, and other neurodegenerative conditions are caused by misfolded prion proteins has had significant implications for our understanding of neurodegenerative disease. Prion diseases are invariably fatal and currently untreatable, and the emergence of bovine spongiform encephalopathy (BSE, or ``mad cow disease'') in the 1980s and 1990s---and its apparent transmission to humans as variant Creutzfeldt-Jakob disease---has made the study of prions a matter of urgent public health concern.

The prion hypothesis has also influenced our understanding of other neurodegenerative diseases, including Alzheimer's disease and Parkinson's disease, in which misfolded proteins appear to play a central role. The concept that protein misaggregation can be self-propagating and transmissible---a concept that originated with Gajdusek's work on kuru---is now a major area of research in neurology and neuroscience.

Moreover, Gajdusek's work on the epidemiology of kuru in Papua New Guinea demonstrated the importance of cultural and behavioral factors in the transmission of infectious diseases---a lesson that remains relevant in an era of emerging and re-emerging infections. The recognition that funerary practices, dietary habits, and other cultural practices can influence the spread of disease has informed public health interventions in many parts of the world, from the control ofEbola virus disease in West Africa to the management of variant Creutzfeldt-Jakob disease in the United Kingdom.

\section{Legacy}

Baruch Blumberg's legacy extends well beyond his Nobel Prize-winning work on hepatitis B. After the discovery of HBV, Blumberg turned his attention to the broader question of the relationship between infectious agents and chronic disease. He founded the National Center for Human Genome Research at the NIH (later renamed the National Human Genome Research Institute) and played a key role in the planning and launch of the Human Genome Project. In his later years, Blumberg became interested in the application of systems biology and genomic approaches to global health, and he founded the Institute for Systems Biology in Seattle, Washington. Blumberg died in 2011, but his influence on genomics, global health, and cancer research continues to be felt.

D. Carleton Gajdusek's legacy is more complex. His work on kuru and prion diseases opened up an entirely new field of biology, and his demonstration that an infectious agent could cause a neurodegenerative disease without involving nucleic acids was one of a significant insights in the history of medicine. However, Gajdusek's personal life was marked by controversy. In the 1990s, he was convicted of child molestation involving boys he had brought from Micronesia to the United States, and he spent time in prison before his death in 2008. Despite these personal failings, his scientific contributions remain of enduring importance.

The work of both laureates illustrates a broader principle: that important discoveries in medicine often come from the study of unusual or neglected diseases in unusual or neglected populations. Blumberg's study of genetic variation in remote populations led him to the Australia antigen and, ultimately, to hepatitis B virus. Gajdusek's investigation of a mysterious disease in a remote tribe in Papua New Guinea led him to the discovery of prions. In both cases, the willingness to venture beyond the familiar and to study disease in its natural setting proved to be the key to unlocking fundamental biological truths.


\section{Modern Developments}

Blumberg and Gajdusek's work has led to modern prion research and hepatitis B control. Blumberg's hepatitis B vaccine (first cancer-preventing vaccine) has dramatically reduced liver cancer worldwide. Understanding of prion diseases has informed research on Alzheimer's and Parkinson's diseases, which may involve prion-like mechanisms. Antiviral drugs for hepatitis B (entecavir, tenofovir) suppress the virus but do not cure it, driving research for functional cures.


\section{Bibliography}

\begin{thebibliography}{9}
\bibitem{nobel-1976}
Nobel Prize Outreach, ``The Nobel Prize in Physiology or Medicine 1976,''\emph{NobelPrize.org}.\\
\url{https://www.nobelprize.org/prizes/medicine/1976/summary/}

\bibitem{lecture-1976}
Baruch S. Blumberg, ``Nobel Lecture,''\emph{NobelPrize.org}. Winner-authored primary source; downloadable from the official Nobel Prize archive.\\
\url{https://www.nobelprize.org/prizes/medicine/1976/blumberg/lecture/}
\end{thebibliography}
